% Section 8.2, from lectures/L08/README.md: the objectives, the evaluation questions, and the next
% lecture.

\section{Review}
\label{c8:sec:review}

\needspace{6\baselineskip}
\subsection*{What you should be able to do now}
After this chapter, you should be able to:
\begin{itemize}
  \item \textbf{Build a scripted stub} and use it to assert the exact transactions the driver
    produces, including the byte order and the ``no \code{RX_POP} on empty'' behaviour.
  \item \textbf{Explain what a test double has to be faithful about} and what it may ignore, and why
    the stub's two buffers are linear logs rather than FIFOs.
  \item \textbf{Read a red test and place the bug}, distinguishing a fault in the stub's script from
    a fault in the driver it is scripting.
  \item \textbf{Build blocking helpers over a non-blocking core}, and explain why the spinning lives
    above the driver rather than inside it.
\end{itemize}

\needspace{6\baselineskip}
\subsection*{Questions to test yourself}
You should be able to explain:
\begin{enumerate}
  \item Why is the driver tested against a scripted \code{driver::transport::Stub} rather than the
    real SPI at this stage, and what class of bug does that let you find first?
  \item \code{injectRxWord} queues its four bytes least significant first by mistake. Which cases go
    red, and why does the failure point at \code{readReg()} rather than at the stub?
  \item The suite counts \code{begin()} calls in eight cases, and \code{end()} calls in two of
    them. What would be wrong
    with a driver that passed every byte-sequence assertion but failed those counters?
  \item Why do \code{writeBlocking} and \code{readBlocking} take \code{Interface&} rather than
    \code{Uart&}, and which class of \chapref{6} does that choice make usable in \chapref{10}?
\end{enumerate}

\needspace{6\baselineskip}
\subsection*{Looking ahead}
In \chapref{9} the bottom layer becomes real: the AVR's own \code{SPCR}, \code{SPSR} and
\code{SPDR}, \code{AvrSpi} implementing \code{driver::transport::Interface}, and what a freestanding
target removes. The driver, the interface, the register map and this suite all run unchanged.
